\documentclass[portrait,a0paper,fontscale=0.292]{baposter}


\usepackage{graphicx}
\usepackage{relsize}
\usepackage{ae}
\usepackage{tabularx}
\usepackage{enumitem}
\usepackage{setspace}
\usepackage{helvet}


\graphicspath{{images/}}

 %%%%%%%%%%%%%%%%%%%%%%%%%%%%%%%%%%%%%%%%%%%%%%%%%%%%%%%%%%%%%%%%%%%%%%%%%%%%%%%%
 %%%% Some math symbols used in the text
 %%%%%%%%%%%%%%%%%%%%%%%%%%%%%%%%%%%%%%%%%%%%%%%%%%%%%%%%%%%%%%%%%%%%%%%%%%%%%%%%



 %%%%%%%%%%%%%%%%%%%%%%%%%%%%%%%%%%%%%%%%%%%%%%%%%%%%%%%%%%%%%%%%%%%%%%%%%%%%%%%%
 % Define Unibas Colors
 %%%%%%%%%%%%%%%%%%%%%%%%%%%%%%%%%%%%%%%%%%%%%%%%%%%%%%%%%%%%%%%%%%%%%%%%%%%%%%%%
 \definecolor{color1}{RGB}{165,215,210}
 \definecolor{color2}{RGB}{30,165,165}
 \definecolor{color3}{RGB}{45,55,60}
 \definecolor{color4}{RGB}{140,145,150}
 \definecolor{color5}{RGB}{210,5,55}
 \definecolor{color6}{RGB}{235,130,155}
 \definecolor{color7}{RGB}{0,110,110}
 \definecolor{color8}{RGB}{190,195,200}
 \definecolor{color9}{RGB}{0,0,0}
 \definecolor{color10}{RGB}{255,255,255}

 %%%%%%%%%%%%%%%%%%%%%%%%%%%%%%%%%%%%%%%%%%%%%%%%%%%%%%%%%%%%%%%%%%%%%%%%%%%%%%%%
 % Multicol Settings
 %%%%%%%%%%%%%%%%%%%%%%%%%%%%%%%%%%%%%%%%%%%%%%%%%%%%%%%%%%%%%%%%%%%%%%%%%%%%%%%%
 \setlength{\columnsep}{0.7em}
 \setlength{\columnseprule}{0mm}


 %%%%%%%%%%%%%%%%%%%%%%%%%%%%%%%%%%%%%%%%%%%%%%%%%%%%%%%%%%%%%%%%%%%%%%%%%%%%%%%%
 % Save space in lists. Use this after the opening of the list
 %%%%%%%%%%%%%%%%%%%%%%%%%%%%%%%%%%%%%%%%%%%%%%%%%%%%%%%%%%%%%%%%%%%%%%%%%%%%%%%%
 \newcommand{\compresslist}{%
 \setlength{\itemsep}{1pt}%
 \setlength{\parskip}{0pt}%
 \setlength{\parsep}{0pt}%
 }
 
 %%%%%%%%%%%%%%%%%%%%%%%%
 %\renewcommand*{\familydefault}{\sfdefault}

%-----------------------------------------------------------------------------

%%%%%%%%%%%%%%%%%%%%%%%%%%%%%%%%%%%%%%%%%%%%%%%%%%%%%%%%%%%%%%%%%%%%%%%%%%%%%
%% Begin of Document
%%%%%%%%%%%%%%%%%%%%%%%%%%%%%%%%%%%%%%%%%%%%%%%%%%%%%%%%%%%%%%%%%%%%%%%%%%%%%
\begin{document}
%%%%%%%%%%%%%%%%%%%%%%%%%%%%%%%%%%%%%%%%%%%%%%%%%%%%%%%%%%%%%%%%%%%%%%%%%%%%%
%% Here starts the poster
%%---------------------------------------------------------------------------
%% Format it to your taste with the options
%%%%%%%%%%%%%%%%%%%%%%%%%%%%%%%%%%%%%%%%%%%%%%%%%%%%%%%%%%%%%%%%%%%%%%%%%%%%%
\begin{poster}{
 % Show grid to help with alignment
 grid=false,
 % Column spacing
 colspacing=0.7em,
 % Color style
 headerColorOne=color1,
 headerColorTwo=color1,
 borderColor=color1,
 headerfont=\Large\bf\textsf ,
 textfont=\sf,
 % Format of textbox
 textborder=faded,
 boxColorOne=color10,
 % Format of text header
 headerborder=open,
 headershape=roundedright,
 headershade=plain,
 background=none,
 %bgColorOne=color10,
 %bgColorTwo=color10,
 headerheight=0.12\textheight}
 % Eye Catcher
 {
	\parbox[t]{0.2\textwidth}{
	\vspace{-2.5\baselineskip}
 	\includegraphics[height=0.05\textheight]{logo}
 	}
 }
 % Title
 {
	\sf\textbf{Towards an Exascale-Ready Mini-App for Smooth Particle Hydrodynamics}
 }
 % Authors
 {
 	\sf
 	\vspace{0.5\baselineskip}
 	%\textbf{Danilo Guerrera}, Helmar Burkhart, and Antonio Maffia\\[1em]
 	Florina M. Ciorba, Lucio Mayer, Rubén Cabezon,David Imbert,\\
 	 Danilo Guerrera, Aurélien Cavelan, Darren S. Reed, Jean-Guillaume Piccinali,\\
 	 Ioana Banicescu, Domingo Garciá-Senz, and Thomas R. Quinn\\[1em]
 	%\texttt{danilo.guerrera@unibas.ch}
 	%\texttt{danilo.guerrera@unibas.ch}
 }
 % University logo
 {
 %\includegraphics[height=0.05\textheight]{logo}
 }

%%%%%%%%%%%%%%%%%%%%%%%%%%%%%%%%%%%%%%%%%%%%%%%%%%%%%%%%%%%%%%%%%%%%%%%%%%%%%%
%%% Now define the boxes that make up the poster
%%%---------------------------------------------------------------------------
%%% Each box has a name and can be placed absolutely or relatively.
%%% The only inconvenience is that you can only specify a relative position 
%%% towards an already declared box. So if you have a box attached to the 
%%% bottom, one to the top and a third one which should be inbetween, you 
%%% have to specify the top and bottom boxes before you specify the middle 
%%% box.
%%%%%%%%%%%%%%%%%%%%%%%%%%%%%%%%%%%%%%%%%%%%%%%%%%%%%%%%%%%%%%%%%%%%%%%%%%%%%%

%%%%%%%%%%%%%%%%%%%%%%%%%%%%%%%%%%%%%%%%%%%%%%%%%%%%%%%%%%%%%%%%%%%%%%%%%%%%%%
	\headerbox{1. Vision}{name=objectives,column=0, row=0, span=1}{
%%%%%%%%%%%%%%%%%%%%%%%%%%%%%%%%%%%%%%%%%%%%%%%%%%%%%%%%%%%%%%%%%%%%%%%%%%%%%%
		\smallskip
	
%		\begin{quote}``Non-reproducible single occurrences are of no 
%		significance to science.'' \textit{(Karl Popper, 1959)}\end{quote}

		\begin{itemize}[leftmargin=0.08\textwidth]
		%\compresslist
			\item Evaluate Smooth Particle Hydrodynamics codes in terms of performance
			\item Understand the limits of the actual implementations
			\item Derive a mini-app that synthesizes their characteristics
			\item Provide a optimized implementations of basics SPH operands
		\end{itemize}
		
   }
%%%%%%%%%%%%%%%%%%%%%%%%%%%%%%%%%%%%%%%%%%%%%%%%%%%%%%%%%%%%%%%%%%%%%%%%%%%%%%



%%%%%%%%%%%%%%%%%%%%%%%%%%%%%%%%%%%%%%%%%%%%%%%%%%%%%%%%%%%%%%%%%%%%%%%%%%%%%%
	\headerbox{2. Smooth Particle Hydrodynamics codes}{name=taxonomy,column=1,row=0, span=2}{
%%%%%%%%%%%%%%%%%%%%%%%%%%%%%%%%%%%%%%%%%%%%%%%%%%%%%%%%%%%%%%%%%%%%%%%%%%%%%%
		\begin{minipage}[t]{0.5\textwidth}
			\includegraphics[width=1.2\textwidth]{laplacian}
		\end{minipage} 
		\hfill
		\begin{minipage}[b]{0.4\textwidth}
			\begin{itemize}
				\item purely Lagrangian method (meshless)
				\item used in numerical simulations of fluids in astrophysics and computational fluid dynamics
				\item computationally demanding 
			\end{itemize} 
		\end{minipage}
		
		\textbf{SPH Characteristics}
		\vspace{-1mm}
		\begin{itemize}
			\item purely Lagrangian method (meshless)
			\item used in numerical simulations of fluids in astrophysics and computational fluid dynamics
			\item computationally demanding 
		\end{itemize}
	}
%%%%%%%%%%%%%%%%%%%%%%%%%%%%%%%%%%%%%%%%%%%%%%%%%%%%%%%%%%%%%%%%%%%%%%%%%%%%%%

%%%%%%%%%%%%%%%%%%%%%%%%%%%%%%%%%%%%%%%%%%%%%%%%%%%%%%%%%%%%%%%%%%%%%%%%%%%%%%
	\headerbox{3. Rotating Square Patch}{name=stencil,column=0,below=taxonomy, span=1}{
%%%%%%%%%%%%%%%%%%%%%%%%%%%%%%%%%%%%%%%%%%%%%%%%%%%%%%%%%%%%%%%%%%%%%%%%%%%%%%

		The stencil motif has manifold applications in science, ranging from 
		weather forecast to image processing.
		\begin{center}
			\begin{minipage}[t]{0.35\textwidth}
				\includegraphics[width=\textwidth]{laplacian}
				\centering\hspace*{1ex}\hspace{1ex}\scriptsize 13 Points Stencil
			\end{minipage}
		\end{center}	
		Stencil computations are characterized by low arithmetic intensity: 
		they are memory-bound.
		\begin{center}
			\begin{minipage}[t]{1\textwidth}
				\includegraphics[width=\textwidth]{ai.pdf}
				\centering\hspace*{1ex}\hspace{1ex}\hspace{1ex}\scriptsize Arithmetic Intensity
			\end{minipage}
		\end{center}	
		Sophisticated parallelization techniques are required in order to arrive 
		at scalable solutions: several stencil compilers are available.
		%several stencil compilers have been implemented.
   }
%%%%%%%%%%%%%%%%%%%%%%%%%%%%%%%%%%%%%%%%%%%%%%%%%%%%%%%%%%%%%%%%%%%%%%%%%%%%%%

%%%%%%%%%%%%%%%%%%%%%%%%%%%%%%%%%%%%%%%%%%%%%%%%%%%%%%%%%%%%%%%%%%%%%%%%%%%%%%
	\headerbox{4. Experimental results}{name=experiment,column=1,below=taxonomy, span=2}{
%%%%%%%%%%%%%%%%%%%%%%%%%%%%%%%%%%%%%%%%%%%%%%%%%%%%%%%%%%%%%%%%%%%%%%%%%%%%%%

		\begin{center}
			\begin{minipage}{0.49\textwidth}
			We use two stencil compilers:
				\begin{itemize}
					\compresslist
					\item {\Large\textsc{patus}} \cite{PATUS} defines a DSL
					to express stencil and exploits auto-tuning
					\item {\Large\textsc{pluto}} \cite{PLUTO}, a source to source compiler that uses the polyhedral 
					model approach for compiler optimization
				\end{itemize}				 
				Our reference is a naive implementation of the classical 3-D wave equation, with NUMA-aware initialization, 
				parallelized with OpenMP.
			\end{minipage}
			\begin{minipage}{0.49\textwidth}
				\begin{minipage}{0.49\textwidth}
					\includegraphics[width=\textwidth]{laplacian.pdf}
					\centering\hspace*{1ex}\hspace{1ex}\scriptsize Mint, University of Basel
				\end{minipage}
				\begin{minipage}{0.49\textwidth}
					\includegraphics[width=\textwidth]{laplacian.pdf}
					\centering\hspace*{1ex}\hspace{1ex}\hspace{1ex}\scriptsize LiMa, University of Erlangen
				\end{minipage}
			\end{minipage}
			\medskip
					
			\begin{minipage}[t]{0.45\textwidth}
				University of Basel 2016\hfill
				\smallskip
				{\tiny\setlength\tabcolsep{1.5pt}
				\begin{tabularx}{\textwidth}{|>{\centering}m{0.16\textwidth}|m{0.806\textwidth}|}
				\hline
				\scriptsize\textbf{Problem} &
				\vspace{1mm}
				Calculate a 3-D wave equation of $200^3$ elements (IEEE single precision arithmetic) in 100 timesteps\vspace{0.5mm}	 \\
				\hline
				\scriptsize\textbf{System} &
				\vspace{-2.5mm}
				\begin{description}[labelsep=0.5mm]
					\itemsep0pt \parskip0pt \parsep0pt %\compresslist
					\item[SW:] OpenMP 4.0, GCC 4.9.2, PATUS 0.1.4, PLUTO 0.10
					\item[HW:] 1 node
					\vspace{-1mm}%\vspace{-3.5mm}%\vspace{-0.05\textwidth}
					\begin{itemize}[leftmargin=-0.02\textwidth, labelsep=0.5mm]
						\renewcommand{\labelitemi}{\boldmath\scriptsize$\cdot$}
						\itemsep0pt \parskip0pt \parsep0pt %\compresslist 
						\item CPU: 2x AMD Opteron 6274 "Bulldozer" 16-Core, 2.2 GHz, 12 MiB L3 cache, 4 NUMA domains
						\item RAM: 256 GiB
						\item OS: Ubuntu 14.04.4, Kernel 3.8.0-38
					\end{itemize}
					\vspace{-3.5mm} %\vspace*{-\baselineskip}
				\end{description} \\
				\hline
				\scriptsize\textbf{Method} &
				%\vspace{-0.5mm}
				\begin{enumerate}[leftmargin=0.035\textwidth, labelsep=0.5mm]
					\itemsep0pt \parskip0pt \parsep0pt %\compresslist
					\item Naive OpenMP implementation with NUMA aware initialization (16~ FLOPS)
					\item DSL + auto-tuning with PATUS (20 FLOPS)
					\item Polyhedral model with PLUTO (16 FLOPS)
					\vspace{0.15mm}
				\end{enumerate} \\
				\hline
				\end{tabularx}
				}
			\end{minipage}
			\begin{minipage}[t]{0.45\textwidth}
				University of Erlangen 2016\hfill
				\smallskip
				{\tiny\setlength\tabcolsep{1.5pt}
				\begin{tabularx}{\textwidth}{|>{\centering}m{0.16\textwidth}|m{0.806\textwidth}|}
				\hline
				\scriptsize\textbf{Problem} &
				\vspace{1mm}
				Calculate a 3-D wave equation of $200^3$ elements (IEEE single precision arithmetic) in 100 timesteps\vspace{0.5mm}	 \\
				\hline
				\scriptsize\textbf{System} &
				\vspace{-2.5mm}
				\begin{description}[labelsep=0.5mm] 
					\itemsep0pt \parskip0pt \parsep0pt %\compresslist
					\item[SW:] OpenMP 4.0, GCC 4.9.2, PATUS 0.1.4, PLUTO 0.10
					\item[HW:] 1 node
					\vspace{-1mm}%\vspace{-3.5mm}%\vspace{-0.05\textwidth}
					\begin{itemize}[leftmargin=-0.02\textwidth, labelsep=0.5mm]
						\renewcommand{\labelitemi}{\boldmath\scriptsize$\cdot$}
						\itemsep0pt \parskip0pt \parsep0pt %\compresslist 
						\item CPU: 2x Xeon 5650 "Westmere" 6 cores + SMT, 2.66 GHz, 12 MiB Shared Cache per chip, 2 NUMA domains
						\item RAM: 24 GB (DDR3-1333)
						\item OS: CentOS 6.7, Kernel 2.6.32-573.7.1.el6%.x86\_64
					\end{itemize}
					\vspace{-3.5mm} %\vspace*{-\baselineskip}
				\end{description} \\
				\hline
				\scriptsize\textbf{Method} &
				%\vspace{-0.5mm}
				\begin{enumerate}[leftmargin=0.035\textwidth, labelsep=0.5mm]
					\itemsep0pt \parskip0pt \parsep0pt %\compresslist
					\item Naive OpenMP implementation with NUMA aware initialization (16~ FLOPS)
					\item DSL + auto-tuning with PATUS (20 FLOPS)
					\item Polyhedral model with PLUTO (16 FLOPS)
					\vspace{0.15mm}
				\end{enumerate} \\
				\hline
				\end{tabularx}
				}
			\end{minipage}
						\begin{minipage}[t]{0.45\textwidth}
				University of Erlangen 2016\hfill
				\smallskip
				{\tiny\setlength\tabcolsep{1.5pt}
				\begin{tabularx}{\textwidth}{|>{\centering}m{0.16\textwidth}|m{0.806\textwidth}|}
				\hline
				\scriptsize\textbf{Problem} &
				\vspace{1mm}
				Calculate a 3-D wave equation of $200^3$ elements (IEEE single precision arithmetic) in 100 timesteps\vspace{0.5mm}	 \\
				\hline
				\scriptsize\textbf{System} &
				\vspace{-2.5mm}
				\begin{description}[labelsep=0.5mm] 
					\itemsep0pt \parskip0pt \parsep0pt %\compresslist
					\item[SW:] OpenMP 4.0, GCC 4.9.2, PATUS 0.1.4, PLUTO 0.10
					\item[HW:] 1 node
					\vspace{-1mm}%\vspace{-3.5mm}%\vspace{-0.05\textwidth}
					\begin{itemize}[leftmargin=-0.02\textwidth, labelsep=0.5mm]
						\renewcommand{\labelitemi}{\boldmath\scriptsize$\cdot$}
						\itemsep0pt \parskip0pt \parsep0pt %\compresslist 
						\item CPU: 2x Xeon 5650 "Westmere" 6 cores + SMT, 2.66 GHz, 12 MiB Shared Cache per chip, 2 NUMA domains
						\item RAM: 24 GB (DDR3-1333)
						\item OS: CentOS 6.7, Kernel 2.6.32-573.7.1.el6%.x86\_64
					\end{itemize}
					\vspace{-3.5mm} %\vspace*{-\baselineskip}
				\end{description} \\
				\hline
				\scriptsize\textbf{Method} &
				%\vspace{-0.5mm}
				\begin{enumerate}[leftmargin=0.035\textwidth, labelsep=0.5mm]
					\itemsep0pt \parskip0pt \parsep0pt %\compresslist
					\item Naive OpenMP implementation with NUMA aware initialization (16~ FLOPS)
					\item DSL + auto-tuning with PATUS (20 FLOPS)
					\item Polyhedral model with PLUTO (16 FLOPS)
					\vspace{0.15mm}
				\end{enumerate} \\
				\hline
				\end{tabularx}
				}
			\end{minipage}
		\end{center}	

%		\begin{center}
			
%		\end{center}	

   }
%%%%%%%%%%%%%%%%%%%%%%%%%%%%%%%%%%%%%%%%%%%%%%%%%%%%%%%%%%%%%%%%%%%%%%%%%%%%%%

%%%%%%%%%%%%%%%%%%%%%%%%%%%%%%%%%%%%%%%%%%%%%%%%%%%%%%%%%%%%%%%%%%%%%%%%%%%%%%
%	\headerbox{5. \textsc{prova!}: Architecture Overview}{name=architecture,column=1,below=stencil, span=2}{
%%%%%%%%%%%%%%%%%%%%%%%%%%%%%%%%%%%%%%%%%%%%%%%%%%%%%%%%%%%%%%%%%%%%%%%%%%%%%%
%		\textsc{prova!} is the tool we develop, aiming to manage software stack 
%		(dependencies and environment), build process, running and post-processing 
%		phase, as well as the raw data for generating the performance graphs.
%		to provide an easy way to configure and run an experiment without 
%		dealing with the execution environment and then visualize the results 
%		in a meaningful way. Through a web interface, an 	user can access the 
%		framework, thus interacting with the parallel machines.
%		\begin{center}
%			\includegraphics[width=0.7\linewidth]{prova_arch_daniloX3}
%		\end{center}
%		The tool manages: software stack (dependencies and environment), build 
%		process, running and post-processing phase, as well as the raw data 
%		for generating the graphs.
		
%	}
%%%%%%%%%%%%%%%%%%%%%%%%%%%%%%%%%%%%%%%%%%%%%%%%%%%%%%%%%%%%%%%%%%%%%%%%%%%%%%

 
%%%%%%%%%%%%%%%%%%%%%%%%%%%%%%%%%%%%%%%%%%%%%%%%%%%%%%%%%%%%%%%%%%%%%%%%%%%%%%
	\headerbox{5. SPH-EXA Mini-app}{name=pillars,column=0,below=stencil, span=1}{
%%%%%%%%%%%%%%%%%%%%%%%%%%%%%%%%%%%%%%%%%%%%%%%%%%%%%%%%%%%%%%%%%%%%%%%%%%%%%%

		\bigskip
		
		What we have:
		\begin{itemize}
			\item extrapolate common features
			\item provide a reference optimized implementation (MPI+X)
			\item provide a library for SPH simulations
		\end{itemize}
		
		\medskip
		
		What we pursue:
		\begin{itemize}
		\compresslist
			\item Integration of likwid \cite{Wellein} into \textsc{prova!}
			\item Evaluation of new compilers
			\item Stencil applications tuning
			\item Develop a performance model
		\end{itemize}
		
		\medskip
		
		Interested?
		\begin{itemize}
		\compresslist
		\item https://prova.io
		\end{itemize}

   }
%%%%%%%%%%%%%%%%%%%%%%%%%%%%%%%%%%%%%%%%%%%%%%%%%%%%%%%%%%%%%%%%%%%%%%%%%%%%%%



%%%%%%%%%%%%%%%%%%%%%%%%%%%%%%%%%%%%%%%%%%%%%%%%%%%%%%%%%%%%%%%%%%%%%%%%%%%%%%
	\headerbox{6. References}{name=references,column=0,below=pillars, span=3}{
%%%%%%%%%%%%%%%%%%%%%%%%%%%%%%%%%%%%%%%%%%%%%%%%%%%%%%%%%%%%%%%%%%%%%%%%%%%%%%
		\smaller
     	\smaller
		\bibliographystyle{plain}
		\renewcommand{\section}[2]{\vskip 0.05em}
		\begin{thebibliography}{1}
			\itemsep=-0.01em
			\setlength{\baselineskip}{0.4em}
			
			\bibitem{Maffia}
			Antonio Maffia, Helmar Burkhart and Danilo Guerrera.
			\newblock Reproducibility in Practice: Lessons Learned from Research 
			and Teaching Experiments.
			\newblock In {\em Euro-Par 2015: Parallel Processing Workshops}, 2015.
			
			\bibitem{Wellein}
			Jan Treibig, Georg Hager, and Gerhard Wellein.
			\newblock LIKWID: A lightweight performance-oriented tool suite for 
			x86 multicore environments.
			\newblock In {\em Proceedings of PSTI2010, the First International 
			Workshop on Parallel Software Tools and Tool Infrastructures}, 2010.
			
			\bibitem{PATUS}
			Matthias Christen, Olaf Schenk, and Helmar Burkhart.
			\newblock PATUS: A Code Generation and Autotuning Framework for Parallel 
			Iterative Stencil Computations on Modern Microarchitectures.
			\newblock In {\em Proceedings of the 2011 IEEE International Parallel 
			\& Distributed Processing Symposium}, 2011.
			
			\bibitem{PLUTO}
			Uday Bondhugula, Albert Hartono, J. Ramanujam, and P. Sadayappan.
			\newblock A practical automatic polyhedral parallelizer and locality 
			optimizer.
			\newblock In {\em Proceedings of the 29th ACM SIGPLAN Conference on 
			Programming Language Design and Implementation}, 2008.
			
		\end{thebibliography}
	}
%%%%%%%%%%%%%%%%%%%%%%%%%%%%%%%%%%%%%%%%%%%%%%%%%%%%%%%%%%%%%%%%%%%%%%%%%%%%%%

\end{poster}%
%
\end{document}
